\begin{table}[!htbp]
\centering\small
\caption{A Joint-Region Certificate for the Adjustment Effect}
\label{tab:r7_decision}
\begin{tabular}{lrrr}
\toprule
Adjustment regime & Lower $\Delta$ & Upper $\Delta$ & Max. class gap \\
\midrule
Adjustment available & $6.88873\times10^{-5}$ & 0.000324537 & $7.5953\times10^{-9}$ \\
No deliberate adjustment & -0.000369485 & $-5.29937\times10^{-5}$ & $2.19355\times10^{-6}$ \\
\bottomrule
\end{tabular}
\par\smallskip
\begin{minipage}{0.98\linewidth}\footnotesize
Notes: Initial state $(u,X)=(2,1.25)$, $k=2$, full finite target, entire $(\lambda,d)\in[0,0.25]\times[0.4,0.45]$. Bounds include the arithmetic allowance. Maximum class gap is the larger of the two class-specific upper-minus-feasible-lower certificates in that regime. The per-class arithmetic allowance is $3.79765\times10^{-9}$; the lower bound on the relative adjustment option is 0.000121881. Eight corner optimizations produce sixteen feasible class policies; eight count constructions certify the rectangle. No generic $10^{-3}$ welfare tolerance is used for these signs.
\end{minipage}
\end{table}
